\documentclass[11pt,a4paper]{article}
% ============================================================
% Shared preamble for MCA101 exam-preparation notes.
%
% libertinus and tcolorbox are installed in the user TeX tree
% (/home/elish4h/texmf) via:
%   tlmgr --usermode install --ignore-warning \
%     tcolorbox libertinus libertinus-type1 libertinus-otf \
%     libertinus-fonts libertinust1math
%
% The box styling matches the repository's notes-template.tex.
% ============================================================

\usepackage[margin=0.72in]{geometry}
\usepackage{libertinus}
\usepackage{microtype}
\usepackage[T1]{fontenc}
\usepackage{xcolor}
\usepackage{enumitem}
\usepackage{tabularx}
\usepackage{array}
\usepackage{booktabs}
\usepackage{longtable}
\usepackage{amsmath}
\usepackage{amssymb}
\usepackage{fancyhdr}
\usepackage{titlesec}
\usepackage[skins,breakable]{tcolorbox}
\usepackage{hyperref}
\usepackage{multicol}
\usepackage{parskip}

\definecolor{navy}{HTML}{183153}
\definecolor{blue}{HTML}{2563EB}
\definecolor{green}{HTML}{166534}
\definecolor{red}{HTML}{B91C1C}
\definecolor{gold}{HTML}{92400E}
\definecolor{lightblue}{HTML}{EFF6FF}
\definecolor{lightgreen}{HTML}{F0FDF4}
\definecolor{lightred}{HTML}{FEF2F2}
\definecolor{lightgold}{HTML}{FFFBEB}
\definecolor{grey}{HTML}{475569}

\hypersetup{colorlinks=true,linkcolor=blue,urlcolor=blue}

\titleformat{\section}{\Large\bfseries\color{navy}}{}{0pt}{}
\titleformat{\subsection}{\large\bfseries\color{navy}}{}{0pt}{}
\titleformat{\subsubsection}{\normalsize\bfseries\color{blue}}{}{0pt}{}

\setlist[itemize]{leftmargin=*,nosep}
\setlist[enumerate]{leftmargin=*,nosep}

% Allow tiny line-length overruns to be absorbed rather than overfull.
\setlength{\emergencystretch}{2em}

\tcbset{
  enhanced,
  breakable,
  colbacktitle=blue,
  coltitle=white,
  fonttitle=\bfseries,
  boxrule=0.5pt,
  arc=1.5pt,
  left=6pt,right=6pt,top=4pt,bottom=4pt,
  before skip=6pt,after skip=6pt,
}

\newtcolorbox{exambox}{colback=lightblue,colframe=blue,title=Exam Focus}
\newtcolorbox{rememberbox}{colback=lightgold,colframe=gold,title=Remember}
\newtcolorbox{dangerbox}{colback=lightred,colframe=red,title=Do Not Confuse}
\newtcolorbox{answerbox}{colback=lightgreen,colframe=green,title=Answer Skeleton}
\newtcolorbox{pyqbox}{colback=white,colframe=grey,title=Past-paper question}
\newtcolorbox{solbox}{colback=lightgreen,colframe=green,title=Solution}

\newcommand{\mkmarks}[1]{\hfill\textbf{[#1 marks]}}
\newcommand{\cmd}[1]{\texttt{#1}}
\newcommand{\F}[1]{\ensuremath{#1}}
\newcommand{\subject}[1]{\lhead{#1}\rhead{MCA101 Exam Notes}\cfoot{\thepage}}

\subject{MCA101 -- Overview and PYQ Analysis}

\title{\textbf{Probability, Statistics and Computational Mathematics}\\}
\author{\large MCA101 Exam Preparation -- Overview, Past-Paper Analysis\\and Study Plan}
\date{Exam: Monday 16 November 2026, 9:30 a.m. to 12:30 p.m.}

\begin{document}
\maketitle

\begin{exambox}
\textbf{How to use this set.} Five module PDFs plus this overview sit in \texttt{Statistics/Exam-Prep/}. Print the module you are revising, work the self-test at the end of each one, then use Section~5 here for the six-day plan and Section~6 as your single-page formula sheet on exam day. Every worked answer in these notes was computed and checked numerically.
\end{exambox}

\tableofcontents
\newpage

\section{What the exam looks like}

The timetable dated 26-10-2026 places \textbf{MCA101 Probability, Statistics and Computational Mathematics} on \textbf{Monday 16-11-2026, 9:30 a.m.\ to 12:30 p.m.} (3 hours). The syllabus has five modules of 12 hours each.

\subsection{Two paper patterns appear in the past papers}

The two most recent full papers do not use the same layout, so do not assume one.

\begin{center}\small
\begin{tabularx}{\textwidth}{>{\bfseries}p{0.18\textwidth} X X}
\toprule Pattern & Nov 2024 paper & Oct 2025 paper\\\midrule
Total marks & 75 & 75\\
Duration & 3 hours & 3 hours\\
Section A & 12 short questions, 3 marks each, answer any TEN (30 marks) & 5 long questions, 15 marks each, answer ALL (75 marks)\\
Section B & 5 long questions, 9 marks each, answer ALL (45 marks) & merged into Section A\\
Practical split & long questions carry an OR pair & long questions carry an OR pair\\\bottomrule
\end{tabularx}
\end{center}

\begin{dangerbox}
Do not over-plan for a single layout. If the paper uses the Nov 2024 shape you must be able to write \textbf{12 quick answers} under time pressure; if it uses the Oct 2025 shape you must be able to write \textbf{5 full derivations}. Practise both: the self-test sections in each module PDF cover both the 3-mark definition style and the 9-to-15-mark problem style.
\end{dangerbox}

\subsection{The two continuous assessment papers (less weight, still indicative)}

\begin{center}\small
\begin{tabularx}{\textwidth}{>{\bfseries}p{0.16\textwidth} p{0.12\textwidth} X}
\toprule Paper & Marks & Focus\\\midrule
CAE-1, Sept 2024 & 40 & Module 1 only, with Bayes' theorem and conditional probability proofs\\
CAE-2, Oct 2024 & 20 & Mixed: set counting, arrangement of letters, Poisson, logic, induction\\
CAE-1, Aug 2025 & 40 & Mixed: Bayes, binomial, logic proofs, rules of inference\\
CAE-2, Oct 2025 & -- & Mixed: same distribution as CAE-1, different numbers\\\bottomrule
\end{tabularx}
\end{center}

\section{The five modules and where the marks come from}

Each module is a separate PDF named \texttt{Module-N-<Topic>.pdf}. All six are in the same folder as this document.

\begin{center}\small
\begin{tabularx}{\textwidth}{>{\bfseries}p{0.10\textwidth} X}
\toprule Module & Content\\\midrule
1 & Probability theory: sample space, events, approaches, addition and multiplication theorems, conditional probability, Bayes' theorem\\
2 & Random variables and distributions: PMF/PDF, marginal and conditional functions, moments and MGF, binomial, Poisson, uniform, normal\\
3 & Basic statistics: central tendency, dispersion, moments, skewness, kurtosis, Karl Pearson and rank correlation, regression\\
4 & Mathematical logic: propositional and predicate logic, equivalences, normal forms, quantifiers, rules of inference\\
5 & Counting and mathematical induction: addition and multiplication principles, permutations, combinations, pigeonhole, inclusion-exclusion, induction\\\bottomrule
\end{tabularx}
\end{center}

\section{Past-paper frequency analysis}

Six MCA101 papers were extracted from \texttt{Statistics/STATI PYQ.pdf} (Sept 2024 CAE-1, Oct 2024 CAE-2, Aug 2025 CAE-1, Nov 2024 full, Oct 2025 full) and analysed topic by topic. Frequency of appearance:

\begin{center}\small
\begin{tabularx}{\textwidth}{>{\bfseries}p{0.30\textwidth} c X}
\toprule Topic & Papers hit & What to expect\\\midrule
Bayes' theorem (state, prove, apply) & 4 of 5 & A near-certain question. It appears as a 9-mark proof and as a 3-mark \emph{state and prove} in both CAE-1 papers.\\
Conditional probability and independence & 4 of 5 & Always a dice, card or urn problem. \emph{Given that} questions dominate.\\
Binomial distribution & 4 of 5 & Both directions: compute a probability, and work backwards to find $p$ from two given probabilities.\\
Logical equivalence, truth tables, tautology proofs & 4 of 5 & Two styles: truth-table proof and algebraic proof without a truth table.\\
Rules of inference and validity of arguments & 3 of 5 & Show that a conclusion follows, or that premises are inconsistent.\\
Counting: committees, arrangements of letters, digit numbers & 3 of 5 & These specific word-problems recur almost verbatim across years.\\
Inclusion-exclusion on three sets & 2 of 5 & Both the counting form and the probability form.\\
Mathematical induction & 2 of 5 & Usually a divisibility statement.\\
Rank correlation and regression & 2 of 5 & Module 3 long question, usually as an OR pair.\\
Central tendency and dispersion & 2 of 5 & Definitions plus one computation.\\
Poisson distribution & 1 of 5 & Appeared in Oct 2024 CAE-2 only, but it is in the syllabus.\\
Normal distribution & 1 of 5 & Appeared as a 3-mark standardisation question.\\
Pigeonhole principle & 1 of 5 & Oct 2024 CAE-2 asked for an explanation with an example.\\\bottomrule
\end{tabularx}
\end{center}

\begin{exambox}
\textbf{The single highest-value thing you can do.} Memorise Bayes' theorem so that you can write both the formula and a full proof without thinking, and drill the \emph{probability table} method for applying it. It is the most repeated topic across every paper in the set, and it is the one where a 9-mark answer is fully deterministic once the table is written.
\end{exambox}

\section{Recurring problems, by topic}

These are the exact questions that repeat, with the year they appeared. Knowing them lets you rehearse.

\subsection{Module 1}
\begin{itemize}
\item \textbf{Bayes with two urns.} A bag has 2 white and 3 black; another has 3 white and 2 black. A ball is drawn from one of the bags and found white. Probability it came from the first bag? (Nov 2024 Q13, Oct 2025 Q1) Answer $\tfrac{2}{5}$.
\item \textbf{HIV test.} 90\% detection rate among HIV-positive, 1\% false positive among HIV-negative, 0.1\% prevalence. Probability a person testing positive actually has HIV? (Nov 2024 Q13, Oct 2025 Q1) Answer $\approx 0.0826$.
\item \textbf{Three machines producing bolts.} Output shares 25/35/40\%, defect rates 5/4/2\%. Given a defective bolt, probability it came from machine B? (Aug 2025 Q2) Answer $\approx 0.4058$.
\item \textbf{Lost card from a deck.} A card is lost; two cards drawn from the remaining 51 are both spades. Probability the lost card was a spade? (Aug 2025 Q2) Answer $\tfrac{11}{50}$.
\item \textbf{Three-set inclusion-exclusion in probability.} Programming languages, vehicles, and so on. (Oct 2024 CAE-2, Nov 2024 Q17, Oct 2025 Q5)
\item \textbf{Dice conditional probabilities.} Sum greater than 9 given a specific die showed 5; sum equal to 8 given the other die was below 4. (Sept 2024 Q2, Oct 2025 Q1)
\end{itemize}

\subsection{Module 2}
\begin{itemize}
\item \textbf{Guess at least 6 of 10 true/false answers.} Binomial with $p=\tfrac12$. (Nov 2024 Section A) Answer $\tfrac{386}{1024}\approx 0.377$.
\item \textbf{Find $p$ and $q$ from two binomial probabilities.} Six trials, $P(3)=0.2457$, $P(4)=0.0819$. (Oct 2025 Q2) Solve by dividing the two equations.
\item \textbf{Defective bulbs, sampling with replacement.} 30 bulbs with 6 defective, sample 4 with replacement, find the distribution. (Oct 2025 Q2)
\item \textbf{Expected number of white balls and expected reward.} Bag of 4 white and 6 red, draw 3; then rewards of Rs.4 and Rs.6. (Nov 2024 Q14, Oct 2025 Q2) Answers $1.2$ and Rs.\ $15.60$.
\item \textbf{Normal distribution standardisation.} $\mu=50$, $\sigma=8$, find $P(34<X<62)$. (Nov 2024 Section A) Answer $\approx 0.9104$.
\end{itemize}

\subsection{Module 3}
\begin{itemize}
\item \textbf{Three judges rank ten contestants.} Find which pair of judges agrees most, using rank correlation. (Nov 2024 Q15) Judges 1 and 3, with $\rho\approx 0.636$.
\item \textbf{Husband's and wife's age, two regression lines.} (Oct 2025 Q3) $y=-1.739+0.891x$ and $x=11.25+0.75y$; husband's age for wife's age 23 is $28.5$.
\item \textbf{Rank correlation between two subjects.} English and Mathematics marks, ten students. (Oct 2025 Q3) $\rho\approx 0.733$.
\item \textbf{Standard deviation of two observations.} Show that the SD of two observations $a$ and $b$ is half the distance between them. (Oct 2025 Q3)
\item \textbf{Mean correction.} Mean of 10 items was 70, one item of 92 was misread as 29. Correct mean? (Oct 2025 Q3) Answer $76.3$.
\end{itemize}

\subsection{Module 4}
\begin{itemize}
\item \textbf{Truth-table proofs.} $\neg(p\leftrightarrow q)\equiv(p\land\neg q)\lor(q\land\neg p)$; $p\land(q\lor r)\equiv(p\land q)\lor(p\land r)$. (Nov 2024 Q16)
\item \textbf{Algebraic proof without a truth table.} $p\leftrightarrow q\equiv(p\land q)\lor(\neg p\land\neg q)$; also $[(p\lor q)\land\neg p]\to\neg q \equiv q\to p$. (Nov 2024 Q16, Aug 2025 Q3)
\item \textbf{Principal disjunctive normal form of $\neg P\lor Q$.} (Oct 2025 Q4)
\item \textbf{Validity of an argument.} Four premises with names, conclude a statement about D. (Oct 2025 Q4, Aug 2025 Q4)
\item \textbf{Show premises are inconsistent.} $P\to Q$, $P\to R$, $Q\to\neg R$, and $P$. (Aug 2025 Q4)
\item \textbf{Predicate calculus with dragons.} Some, no, every, not every dragon sleeping. (Oct 2025 Q4)
\end{itemize}

\subsection{Module 5}
\begin{itemize}
\item \textbf{Committee with at least 3 men.} 7 men and 6 women, choose 5 with at least 3 men. (Nov 2024 Q17, Oct 2025 Q5) Answer $756$.
\item \textbf{Three-digit numbers from 2,3,5,6,7,9 divisible by 5, no repetition.} (Nov 2024 Q17, Oct 2025 Q5) Answer $20$.
\item \textbf{Arrangements of the letters of a word.} ATTITUDE with no two T adjacent (Oct 2024 CAE-2, answer $2400$), ALLAHABAD (Oct 2025, answer $7560$), LEADING with vowels together (Oct 2025).
\item \textbf{Number induction.} Prove $3n-1$ is a multiple of 2 (Oct 2025); prove $2^{2n}-1$ is divisible by 3 (Oct 2024 CAE-2).
\item \textbf{Pigeonhole principle with an example.} (Oct 2024 CAE-2)
\item \textbf{Multiplication rule of counting.} Stated in both Nov 2024 and Oct 2025.
\end{itemize}

\section{Six-day study plan, 10 to 16 November}

The plan assumes roughly 6 hours of study a day. If you have less, protect Sections 1 to 3 of the schedule first.

\begin{center}\small
\begin{tabularx}{\textwidth}{>{\bfseries}p{0.13\textwidth} p{0.22\textwidth} X p{0.20\textwidth}}
\toprule Day & Module & Work & Check\\\midrule
Tue 10 Nov & 1 & Read \texttt{Module-1} in full. Drill Bayes' theorem proof and the probability-table method. Do the Module 1 self-test. & You can write $P(A\mid B)=\frac{P(B\mid A)P(A)}{P(B)}$ and the proof from memory.\\
Wed 11 Nov & 2 & Read \texttt{Module-2}. Master binomial, Poisson, uniform and normal. Work the Oct 2025 bulb and dart problems twice. & You can state mean, variance and MGF for all four distributions without looking.\\
Thu 12 Nov & 3 & Read \texttt{Module-3}. Practise one full regression problem and one rank correlation problem per hour. & You can compute $r$, both regression lines, and $\rho$ from a raw table.\\
Fri 13 Nov & 4 & Read \texttt{Module-4}. Write out all equivalences and all inference rules on one sheet. Practise truth-table and algebraic proofs. & You can produce either style of proof for any equivalence given.\\
Sat 14 Nov & 5 & Read \texttt{Module-5}. Then sit a full mock: attempt all questions from Section A of the Nov 2024 paper under 40 minutes, then the Section B questions under 100 minutes. & Mark yourself honestly and list every topic you missed.\\
Sun 15 Nov & All & Re-read only the boxes you marked as weak. Write the formula sheet from Section 6 from memory on blank paper, then compare. Sleep early. & A clean sheet with no gaps.\\\bottomrule
\end{tabularx}
\end{center}

\begin{exambox}
\textbf{Exam-morning rules.} Bring two pens, a ruler, and the standard normal table if your department issues one. Write the question number in the margin before you start a long answer. For every numerical answer, show the formula, substitute, then state the result with units. A correct answer with no working often loses half the marks; a formula with no substitution usually gets partial credit.
\end{exambox}

\section{Formula sheet}

\subsection{Probability (Module 1)}
\begin{center}\small
\begin{tabularx}{\textwidth}{>{\bfseries}p{0.30\textwidth} X}
\toprule Item & Formula\\\midrule
Complement & $P(A^c)=1-P(A)$\\
Addition theorem & $P(A\cup B)=P(A)+P(B)-P(A\cap B)$; if $A\cap B=\emptyset$ then $P(A\cup B)=P(A)+P(B)$\\
Conditional probability & $P(A\mid B)=\dfrac{P(A\cap B)}{P(B)}$, $P(B)>0$\\
Multiplication theorem & $P(A\cap B)=P(A)P(B\mid A)=P(B)P(A\mid B)$\\
Independent events & $P(A\cap B)=P(A)P(B)$; equivalently $P(A\mid B)=P(A)$\\
Bayes' theorem & $P(A_i\mid B)=\dfrac{P(B\mid A_i)P(A_i)}{\sum_j P(B\mid A_j)P(A_j)}$\\
Odds in favour $a:b$ & $P(A)=\dfrac{a}{a+b}$\\
Classical approach & $P(A)=\dfrac{\#\text{ favourable outcomes}}{\#\text{ equally likely outcomes}}$\\\bottomrule
\end{tabularx}
\end{center}

\subsection{Random variables and distributions (Module 2)}
\begin{center}\small
\begin{tabularx}{\textwidth}{>{\bfseries}p{0.24\textwidth} X}
\toprule Item & Formula\\\midrule
PMF & $p(x)=P(X=x)$, $\sum_x p(x)=1$\\
PDF & $f(x)\ge 0$, $\int_{-\infty}^{\infty} f(x)\,dx=1$\\
CDF & $F(x)=P(X\le x)$, $f(x)=\dfrac{d}{dx}F(x)$\\
Expectation & $E(X)=\sum x\,p(x)$ (discrete), $\int x f(x)\,dx$ (continuous)\\
Variance & $\operatorname{Var}(X)=E(X^2)-[E(X)]^2=E[(X-\mu)^2]$\\
Binomial $B(n,p)$ & $P(X=k)=\binom nk p^k q^{n-k}$; $E=np$; $\operatorname{Var}=npq$; $M_G(t)=(q+pe^t)^n$\\
Poisson $\operatorname{Po}(\lambda)$ & $P(X=k)=\dfrac{e^{-\lambda}\lambda^k}{k!}$; $E=\lambda$; $\operatorname{Var}=\lambda$; $M_G(t)=e^{\lambda(t-1)}$\\
Uniform $U(a,b)$ & $f(x)=\dfrac{1}{b-a}$; $E=\dfrac{a+b}{2}$; $\operatorname{Var}=\dfrac{(b-a)^2}{12}$\\
Normal $N(\mu,\sigma^2)$ & $f(x)=\dfrac{1}{\sigma\sqrt{2\pi}}e^{-(x-\mu)^2/(2\sigma^2)}$; $E=\mu$; $\operatorname{Var}=\sigma^2$\\
Standardisation & $Z=\dfrac{X-\mu}{\sigma}$, then read areas from the $Z$ table\\
Binomial to Poisson & if $n$ large and $p$ small, use $\lambda=np$\\\bottomrule
\end{tabularx}
\end{center}

\subsection{Descriptive statistics (Module 3)}
\begin{center}\small
\begin{tabularx}{\textwidth}{>{\bfseries}p{0.24\textwidth} X}
\toprule Item & Formula\\\midrule
Arithmetic mean & $\bar X=\dfrac{\sum fX}{\sum f}=\dfrac{\sum X}{N}$\\
Median & middle value, or $\dfrac{Q_3+Q_1}{2}$; position $\dfrac{N+1}{2}$\\
Mode & most frequent value\\
Range & $R=x_{\max}-x_{\min}$\\
Coefficient of range & $\dfrac{R}{\bar X}$\\
Quartile deviation & $QD=\dfrac{Q_3-Q_1}{2}$; coefficient of QD $=\dfrac{QD}{\bar X}$\\
Mean deviation & $MD=\dfrac{\sum f|X-A|}{N}$, $A$ the chosen centre\\
Variance & $\sigma^2=\dfrac{\sum f(X-\bar X)^2}{N}$; shortcut $\sigma^2=\dfrac{\sum fX^2}{N}-\bar X^2$\\
Standard deviation & $\sigma=\sqrt{\dfrac{\sum f(X-\bar X)^2}{N}}$\\
Karl Pearson $r$ & $r=\dfrac{n\sum xy-\sum x\sum y}{\sqrt{(n\sum x^2-(\sum x)^2)(n\sum y^2-(\sum y)^2)}}=\dfrac{\sum dx\,dy}{\sqrt{\sum dx^2\,\sum dy^2}}$\\
Rank correlation & $\rho=1-\dfrac{6\sum d^2}{n(n^2-1)}$, $d$ difference of ranks\\
Regression $y$ on $x$ & $y=a+bx$, $b=\dfrac{\sum dx\,dy}{\sum dx^2}=\dfrac{n\sum xy-\sum x\sum y}{n\sum x^2-(\sum x)^2}$, $a=\bar y-b\bar x$\\
Regression $x$ on $y$ & $x=a'+b'y$, $b'=\dfrac{\sum dx\,dy}{\sum dy^2}$, $a'=\bar x-b'\bar y$\\
Link & $r=\sqrt{bb'}$; the two regression lines meet at $(\bar x,\bar y)$\\
Skewness (Pearson) & $\dfrac{3(\bar X-\text{median})}{S}$ or $\dfrac{\bar X-\text{mode}}{S}$\\
Bowley coefficient & $\dfrac{Q_3+Q_1-2\,\text{median}}{Q_3-Q_1}$\\
Kurtosis & $\beta_2=\dfrac{\mu_4}{\mu_2^2}$; $\beta_2=3$ normal, $<3$ platykurtic, $>3$ leptokurtic\\\bottomrule
\end{tabularx}
\end{center}

\subsection{Mathematical logic (Module 4)}
\begin{center}\small
\begin{tabularx}{\textwidth}{>{\bfseries}p{0.30\textwidth} X}
\toprule Item & Formula\\\midrule
Negation & $\neg p$; $p\to q\equiv\neg p\lor q$; $p\leftrightarrow q\equiv(p\to q)\land(q\to p)$\\
Converse / inverse / contrapositive & of $p\to q$: converse $q\to p$; inverse $\neg p\to\neg q$; contrapositive $\neg q\to\neg p$\\
Quantifier negation & $\neg\forall x\,P(x)\equiv\exists x\,\neg P(x)$; $\neg\exists x\,P(x)\equiv\forall x\,\neg P(x)$\\
Rules of inference & modus ponens ($p\to q,\,p\Rightarrow q$); modus tollens ($p\to q,\,\neg q\Rightarrow\neg p$); hypothetical syllogism; disjunctive syllogism; addition; simplification; conjunction; resolution\\
Normal forms & DNF is OR of AND terms; CNF is AND of OR terms; PDNF is the full set of minterms where the formula is true\\
Truth table size & $2^n$ rows for $n$ variables\\\bottomrule
\end{tabularx}
\end{center}

\subsection{Counting and induction (Module 5)}
\begin{center}\small
\begin{tabularx}{\textwidth}{>{\bfseries}p{0.30\textwidth} X}
\toprule Item & Formula\\\midrule
Addition principle & $a_1+a_2+\cdots+a_k$ for mutually exclusive cases\\
Multiplication principle & $a_1a_2\cdots a_k$ for successive stages\\
Permutation & ${}^nP_r=\dfrac{n!}{(n-r)!}$\\
Combination & ${}^nC_r=\dfrac{n!}{r!(n-r)!}$, ${}^nP_r={}^nC_r\,r!$\\
Identical objects & $\dfrac{n!}{n_1!n_2!\cdots n_k!}$\\
Pigeonhole & $N$ objects in $k$ boxes, $N>k$ forces one box with two; general form $\lceil N/k\rceil$; guarantee $r$ in one box needs $k(r-1)+1$\\
Inclusion-exclusion (3 sets) & $|A\cup B\cup C|=|A|+|B|+|C|-|A\cap B|-|A\cap C|-|B\cap C|+|A\cap B\cap C|$\\
Mathematical induction & base case $n_0$; assume $P(k)$; prove $P(k+1)$; conclude for all $n\ge n_0$\\\bottomrule
\end{tabularx}
\end{center}

\section{Common mistakes that cost marks}

\begin{enumerate}
\item Using $P(A\cup B)=P(A)+P(B)$ when the events can overlap. Subtract $P(A\cap B)$ unless mutual exclusivity is given.
\item Reversing $P(A\mid B)$ and $P(B\mid A)$. They are not interchangeable.
\item Forgetting the prior probability in a Bayes table, or dividing by one row instead of the column total.
\item Treating "at least one" as a sum of positive terms when the complement is easier.
\item Assuming independent means mutually exclusive. They are different properties.
\item Dividing by $n$ instead of $n-1$ in a sample standard deviation, or mixing the two conventions.
\item Confusing the order of $x$ and $y$ in a regression line, then quoting the wrong line for the prediction asked.
\item In rank correlation, failing to replace tied ranks with their average rank.
\item In counting words with repeated letters, dividing by the wrong factorial.
\item Writing "we know" instead of showing the formula and the substitution in a numerical answer.
\end{enumerate}

\vfill
\begin{center}\textit{Every numerical answer in this set was computed and checked. Where a past paper's own data is internally inconsistent, the question is flagged rather than silently answered wrongly.}\end{center}

\end{document}
